%! Transformations of Functions — Unit 1, Lesson 1.12 — Lesson Plan
\documentclass[10pt]{article}
\usepackage{precalculus-article}
\usepackage{precalculus-boxes}
\usepackage{graphicx}
\graphicspath{{images/}}
\newcommand{\TallMath}[1]{$\displaystyle #1\rule[-1.4em]{0pt}{3.2em}$}
\newcommand{\CourseName}{AP Precalculus}
\newcommand{\SchoolYear}{2026--2027}
\newcommand{\MeetingLength}{55 minutes}
\newcommand{\UnitNumberName}{Unit 1: Polynomial and Rational Functions \quad}
\newcommand{\LessonNumberName}{Lesson 1.12: Transformations of Functions}

\begin{document}
\begin{center}
    {\Large\bfseries \CourseName: \SchoolYear \\
    {\small\bfseries \UnitNumberName \LessonNumberName}}
\end{center}

\begin{tcolorbox}[colback=sky, colframe=navy, boxrule=0.9pt, arc=2mm,
    left=2mm, right=2mm, top=1.5mm, bottom=1.5mm]
    \textbf{Primary Objective:} Students will identify, describe, and construct transformations
    of functions using additive transformations (vertical and horizontal translations) and
    multiplicative transformations (vertical and horizontal dilations, including reflections),
    and will determine the domain and range of a transformed function from the parent's domain
    and range.
    \hfill\textit{\small Skills: MP 1.C, MP 3.A}
\end{tcolorbox}

\renewcommand{\arraystretch}{1.6}
\begin{skillbox}[Priority Ideas \& Skills]{goldbox}
\begin{tabularx}{\linewidth}{@{}X X@{}}
\textbf{AP Skills} &
\textbf{Key Understandings} \\
\midrule
\textbf{MP 1.C} --- Construct new functions as transformations of known functions. &
\begin{itemize}[leftmargin=1.2em,itemsep=2pt,topsep=0pt]
  \item $g(x)=f(x)+k$: vertical translation by $k$ units. (EK 1.12.A.1)
  \item $g(x)=f(x+h)$: horizontal translation by $-h$ units. (EK 1.12.A.2)
  \item $g(x)=af(x)$: vertical dilation by $|a|$; $a<0$ also reflects over $x$-axis. (EK 1.12.A.3)
  \item $g(x)=f(bx)$: horizontal dilation by $1/|b|$; $b<0$ also reflects over $y$-axis. (EK 1.12.A.4)
\end{itemize} \\
\textbf{MP 3.A} --- Describe characteristics of transformed functions. &
\begin{itemize}[leftmargin=1.2em,itemsep=2pt,topsep=0pt]
  \item Transformations can be combined: $g(x)=af(x+h)+k$. (EK 1.12.A.5)
  \item The domain and range of the transformed function may differ from the parent. (EK 1.12.A.6)
\end{itemize} \\
\end{tabularx}
\end{skillbox}

\begin{skillbox}[{Vocabulary, Concepts, \& Theorems}]{greenbox}
\renewcommand{\arraystretch}{1.4}
\begin{tabularx}{\linewidth}{@{} >{\bfseries\color{navy}}p{0.26\linewidth} X @{}}
parent function         & The original, un-transformed function from which transformed versions are built. \\
transformation          & An operation that shifts, scales, or reflects the graph of a function. \\
vertical translation    & $g(x)=f(x)+k$: shifts the graph up or down by $|k|$ units. \\
horizontal translation  & $g(x)=f(x+h)$: shifts the graph left (if $h>0$) or right (if $h<0$) by $|h|$ units. \\
vertical dilation       & $g(x)=af(x)$: scales outputs by $|a|$; $a<0$ reflects over $x$-axis. \\
horizontal dilation     & $g(x)=f(bx)$: scales inputs by $1/|b|$; $b<0$ reflects over $y$-axis. \\
reflection              & A special dilation with $|a|=1$ or $|b|=1$ and negative sign; flips the graph. \\
\end{tabularx}
\end{skillbox}

\begin{skillbox}[Activate Prior Knowledge \& Spiral Review]{sky}
    \begin{tabularx}{\linewidth}{@{}X@{}}
        Spiral review (warm-up) covers: evaluating $f(x)$, $f(-2)$, and $f(x+1)$ from a formula
        (Lesson 1.1 function notation); reading domain and range from a table (Lesson 1.1);
        interpreting end behavior of a polynomial graph (Lessons 1.6--1.7, prerequisite for
        understanding how dilations affect overall shape).
    \end{tabularx}
\end{skillbox}

\begin{skillbox}[Hook]{sky}
\textbf{``One function --- four disguises. Can you un-transform it back?''}

Display the table of $f(x)=x^2$ alongside tables for $f(x)+3$, $f(x-2)$, and $2f(x)$ without
labeling them. Ask: ``What did I \textit{do} to $f$? How do the outputs change? What stays the
same?'' Let students discuss for 60 seconds before revealing the rules. Emphasize the goal: given
a rule, predict the table; given two tables, reverse-engineer the rule.
\end{skillbox}

\begin{skillbox}[Lesson]{sky}
\begin{multicols}{2}
\textbf{Part 1 (12 min): Four Transformation Types.}

Walk through each type in the guided notes with the $f(x)=x^2$ example:
\begin{enumerate}[leftmargin=1.2em,itemsep=2pt,topsep=2pt]
  \item \textbf{Vertical shift} $g(x)=f(x)+k$: output changes, input doesn't. Domain unchanged; range shifts.
  \item \textbf{Horizontal shift} $g(x)=f(x+h)$: input changes, output values unchanged but at different $x$. Stress the direction reversal: $f(x+3)$ shifts \textit{left}.
  \item \textbf{Vertical dilation} $g(x)=af(x)$: outputs scaled by $a$. Negative $a$ flips.
  \item \textbf{Horizontal dilation} $g(x)=f(bx)$: inputs scaled, graph compressed ($|b|>1$) or stretched.
\end{enumerate}

\columnbreak

\textbf{Part 2 (10 min): Combined Transformations.}

Model $g(x)=-2f(x+1)-3$ step by step:
\begin{enumerate}[leftmargin=1.2em,itemsep=2pt,topsep=2pt]
  \item Replace $x$ with $x+1$ (shift left 1).
  \item Multiply by $-2$ (stretch by 2, reflect over $x$-axis).
  \item Subtract 3 (shift down 3).
\end{enumerate}

Use a table-of-values approach to show how each step changes the outputs numerically.

\textbf{Part 3 (5 min): Domain and Range.}

For $f$ with domain $[-2,4]$ and range $[0,9]$, work out domain/range of $g(x)=-2f(x+1)-3$:
domain shifts left 1 $\to [-3,3]$; range: $[0,9]\times(-2)=[-18,0]$, then $-3 \to [-21,-3]$.
Emphasize that order matters when the factor is negative.
\end{multicols}
\end{skillbox}

\begin{skillbox}[Active Monitoring]{redbox}
\begin{itemize}[leftmargin=1.2em,itemsep=2pt,topsep=0pt]
  \item \textbf{Key misconception to catch:} Students often shift $f(x+3)$ \textit{right} instead of left. Cold-call: ``If $g(x)=f(x+3)$, at what $x$-value does $g$ equal $f(0)$? Why?'' (Answer: $x=-3$, so the graph moved left.)
  \item Circulate during the notes table-fill: check that students compute $g(x)=2f(x)$ by doubling the $f$ column, not the $x$ column.
  \item Watch for students who apply transformations in the wrong order for combined cases.
  \item After Part 3: ask a student to explain why domain shifts left when $h>0$ in $f(x+h)$.
\end{itemize}
\end{skillbox}

\begin{skillbox}[Group Work \& Differentiation]{redbox}
\begin{multicols}{3}
\textbf{Tier R --- Remediate}
\begin{itemize}[leftmargin=1.2em,itemsep=2pt,topsep=2pt]
  \item Pre-factored transformation tables provided.
  \item Match the completed table to the rule (multiple choice).
  \item Focus: reading how outputs change and connecting to the rule label.
\end{itemize}

\columnbreak

\textbf{Tier A --- Apply}
\begin{itemize}[leftmargin=1.2em,itemsep=2pt,topsep=2pt]
  \item Build the table for $g(x)=2f(x-1)+3$ step by step.
  \item State domain and range of $g$.
  \item All computation shown in a scaffolded multi-column table.
\end{itemize}

\columnbreak

\textbf{Tier E --- Extend}
\begin{itemize}[leftmargin=1.2em,itemsep=2pt,topsep=2pt]
  \item Given tables for $f$ and $h$, reverse-engineer the rule for $h$ in terms of $f$.
  \item Identify $a$ and $k$; describe all transformations in words.
\end{itemize}
\end{multicols}
\end{skillbox}

\begin{skillbox}[Individual Work \& Assessment]{redbox}
Exit ticket (3 questions, 1 page):
\begin{enumerate}[leftmargin=1.2em,itemsep=2pt,topsep=0pt]
  \item Identify and describe the transformations in $g(x)=f(x+3)-2$ (tests horizontal shift direction, vertical shift).
  \item Complete the table for $g(x)=-f(x)$ given $f$ values (tests vertical reflection).
  \item State the domain of $g(x)=f(x-4)$ given domain of $f$ is $[-2,6]$ (tests domain shift).
\end{enumerate}
Collect exit tickets; sort by correctness on Q1 (horizontal shift direction) to identify students needing re-teaching before Lesson 1.13.
\end{skillbox}

\begin{skillbox}[Reinforcement \& Extension]{goldbox}
\textbf{Homework:} 5 problems covering all four transformation types (identification and description),
completing transformation tables, domain/range under combined transformations, describing $g(x)=-3f(x+2)+1$
step by step, and an AP-style MC on range under $g(x)=-2f(x)+1$.

\textbf{Reflection:} Students explain in their own words why horizontal transformations act opposite
to the sign of $h$ (builds conceptual understanding for Lesson 1.13).

\textbf{Preview of Lesson 1.13:} Building and selecting models for data --- students will need to
recognize parent function shapes and use transformation parameters to fit a function to a context.
Transformations are the core tool for that lesson.
\end{skillbox}

\end{document}
